\documentclass[10pt,a4paper]{amsart}
\usepackage[margin=19mm]{geometry}
\usepackage{amsmath,amssymb,mathtools}
\usepackage[scr=rsfs,cal=euler]{mathalfa}
\usepackage{xfrac}
\usepackage[unicode,hidelinks,bookmarks=false]{hyperref}
\newcommand{\ie}{i.e.\ }
\newcommand{\cf}{cf.\ }
\newcommand{\resp}{respectively\ }
\newcommand{\Subsection}{\subsection}
\providecommand{\nobreakdash}{\nobreak\hbox{-}}
\setlength{\emergencystretch}{2em}
\multlinegap0pt
\usepackage{mathtools}
\usepackage{amssymb}
\usepackage[shortlabels]{enumitem}
\usepackage{xcolor}
\usepackage{dsfont}
\usepackage{csquotes}
\newtheorem{proposition}{Proposition}
\newcommand{\R}{\mathbb{R}}
\newcommand{\Sect}{\Gamma(\X, \U)}
\newcommand{\U}{\mathcal{U}}
\newcommand{\X}{\mathcal{X}}
\newcommand{\advset}[2]{\mathcal{C}_{#1}(#2)}
\newcommand{\cost}{r}
\newcommand{\costfb}[1]{\cost_{#1}}
\newcommand{\disc}{\gamma}
\newcommand{\f}{f}
\newcommand{\ffb}[1]{\f_{#1}}
\newcommand{\maps}{\text{Map}}
\newcommand{\policy}{\pi}
\newcommand{\state}{x} %
\newcommand{\vvarg}[1]{v_{#1}}
\title{ON REWARD-BALANCING METHODS FOR REINFORCEMENT LEARNING}
\author{Simone Baroncini, Bahman Gharesifard, Giuseppe Notarstefano}
\date{Source: arXiv:2604.20433v1 (CC BY 4.0)}
\begin{document}
\maketitle

\noindent\emph{Source and licence.} ON REWARD-BALANCING METHODS FOR REINFORCEMENT LEARNING. Version arXiv:2604.20433v1 (CC BY 4.0), distributed by arXiv under the Creative Commons Attribution 4.0 International licence, http://creativecommons.org/licenses/by/4.0/, as recorded in the arXiv licence metadata for this exact version and shown on the abstract page https://arxiv.org/abs/2604.20433v1. Full licence notice: \texttt{licenses/arxiv-2604.20433-CC-BY-4.0.txt}.\par\smallskip

\noindent\textit{Editorial scope.} Counted unit: the article's proposition prop:feas\_set\_as\_overestimates (source0.tex lines 885--894). The article's complete original proof of it is delivered with it (source0.tex lines 895--917, one \texttt{proof} environment of 23 lines). No mathematics is written, repaired or re-derived: the statement and the proof are the article's own lines, in the article's own notation, and no retained line is substituted or re-flowed.  Retained uncounted as the source-local context the proof needs: lines 321--323.  Nothing was omitted from any retained span, so the package carries no omission record.  No retained line contains a citation, so the wrapper carries no bibliography and no bibliography entry.  No retained line contains a figure environment, an image-inclusion command or drawing code, so no image is delivered and the package declares no asset dependency.  Editorial formatting only, in addition: an AMS \texttt{amsart} preamble that restates the article's own declarations and macros used by the retained lines, namely the environments \texttt{proposition} on separate counters, together with the standard packages those lines need.  The retained environments print in this document as Proposition 1 in order of appearance, whereas the article prints the same items with the numbering of its own full document.  The complete native source of the article as distributed in the arXiv e-print for this exact version is bundled unchanged as \texttt{sources/198934/source0.tex}.
\begin{align}\label{eq:bellman_equation_v}
    v(\state) = \costfb{\policy}(\state) + \disc  \sum_{z \in \X} \ffb{\policy}(z, \state) v(z), \quad \forall\state \in \X.
\end{align}

\begin{proposition}[\textbf{Admissible inputs as overestimates of the optimal value function}]\label{prop:feas_set_as_overestimates}
    For any nonpositive reward function $\cost$, every function in $\advset{}{\cost}$ is an overestimate of the corresponding optimal value function. Formally:
    \begin{align*}
        \advset{}{\cost} \subseteq \mathcal{C}_*(\cost),
    \end{align*}
    where
    \begin{align}\label{eq:overestimates_set}
        \mathcal{C}_*(\cost) \coloneqq \{\delta \in \maps(\X, \R) : \delta \geq \vvarg{\policy^*\!, \cost}\}.
    \end{align}
\end{proposition}

\begin{proof}
    If $\delta \in \advset{}{\cost}$ then
    \begin{align}\label{eq:inequalities_feasible_set}
        \costfb{\policy}(\state) + (\disc \text{$\ffb{\policy}(\state, \state)$} - 1) \delta(\state) + \sum_{y \neq \state} \disc \ffb{\policy}(y, \state) \delta(y) \leq 0,
    \end{align}
    for every state $\state \in \X$ and any fixed policy $\policy \in \Sect$. Since $\vvarg{\policy\!, \cost}$ satisfies \eqref{eq:inequalities_feasible_set} with equality (by the Bellman equation \eqref{eq:bellman_equation_v}), we can rewrite \eqref{eq:inequalities_feasible_set} as
    \begin{align}\label{eq:error_inequalities}
        \sum_{y \in \X} \disc \ffb{\policy}(y, \state) \tilde{\delta}_{\policy, \cost}(y) \leq \tilde{\delta}_{\policy, \cost}(\state),
    \end{align}
    where we have introduced the shorthand notation $\tilde{\delta}_{\policy, \cost} \coloneqq \delta - \vvarg{\policy\!, \cost}$. Since the left-hand side of \eqref{eq:error_inequalities} is a convex combination scaled by $\disc$, we have
    \begin{align*}
        \disc \min_{z \in \X} \tilde{\delta}_{\policy, \cost}(z) \leq \sum_{y \in \X} \disc \ffb{\policy}(y, \state) \tilde{\delta}_{\policy, \cost}(y) \leq \tilde{\delta}_{\policy, \cost}(\state).
    \end{align*}
    Now, since this holds for all $\state \in \X$, we can conclude that
    \begin{align*}
        \disc \min_{z \in \X} \tilde{\delta}_{\policy, \cost}(z) \leq \min_{x \in \X} \tilde{\delta}_{\policy, \cost}(x),
    \end{align*}
    or
    \begin{align*}
        (1 - \disc)\min_{z \in \X} \tilde{\delta}_{\policy, \cost}(z) \geq 0,
    \end{align*}
    which, since $\disc < 1$, implies that $\tilde{\delta}_{\policy, \cost}(\state) = \delta(\state) - \vvarg{\policy\!, \cost}(\state) \geq 0, \forall \state$. Since the policy $\policy$ is arbitrary, this holds for the optimal one as well, and this concludes the proof.
\end{proof}

\end{document}
